\section{生成模型基础概念}

\subsection{监督学习 vs 无监督学习}

在切入生成模型之前，Justin Johnson 首先梳理了深度学习任务的两个正交维度：

\begin{figure}[H]
\centering
\includegraphics[width=0.95\textwidth]{slides-images/slide-008.jpg}
\caption{监督学习 vs 无监督学习：监督学习有 (X, Y) 数据对，无监督学习只有数据 X\protect\footnotemark}
\end{figure}
\footnotetext{Slides 第9页。}

\textbf{维度一：监督 vs 无监督}
\begin{itemize}
    \item \textbf{监督学习}：有 $(X, Y)$ 数据对，学习从 $X$ 到 $Y$ 的映射（如分类、检测、分割）
    \item \textbf{无监督学习}：只有数据 $X$，从数据本身学习结构（如聚类、PCA、密度估计）
\end{itemize}

\textbf{维度二：判别式 vs 生成式}
\begin{itemize}
    \item \textbf{判别式模型}：学习 $p(Y|X)$，即给定输入 $X$ 预测标签 $Y$ 的概率分布
    \item \textbf{生成式模型}：学习 $p(X)$，即数据本身的概率分布
    \item \textbf{条件生成模型}：学习 $p(X|Y)$，即给定条件 $Y$ 生成数据 $X$ 的概率分布
\end{itemize}

\subsection{概率分布的归一化约束}

理解生成模型的关键在于理解\textbf{概率分布的归一化约束}。

\begin{figure}[H]
\centering
\includegraphics[width=0.95\textwidth]{slides-images/slide-010.jpg}
\caption{判别式模型 vs 生成模型 vs 条件生成模型：三种不同的概率建模方式\protect\footnotemark}
\end{figure}
\footnotetext{Slides 第11页。}

\begin{importantbox}{三种概率模型的本质区别}
\begin{itemize}
    \item \textbf{判别式模型} $p(Y|X)$：对于每个输入 $X$，不同标签 $Y$ 之间竞争概率质量。图片之间没有竞争——无法拒绝不合理的输入。
    \item \textbf{生成模型} $p(X)$：所有可能的图片相互竞争概率质量。模型必须思考什么图片更「合理」。可以通过低概率来拒绝不合理的输入。
    \item \textbf{条件生成模型} $p(X|Y)$：对于每个条件 $Y$，所有可能的图片相互竞争概率质量。这才是最有实用价值的模型。
\end{itemize}
\end{importantbox}

\begin{warningbox}{论文中常省略条件变量 $Y$}
大多数生成模型的论文和教程会写 $p(X)$ 而非 $p(X|Y)$，这是为了简化数学推导。但在实际应用中，几乎所有有实用价值的生成模型都是\textbf{条件生成模型}。阅读论文时应始终意识到：每个 $p(X)$ 背后可能隐含着一个条件变量 $Y$。
\end{warningbox}

\subsection{贝叶斯定理连接三种模型}

三种概率模型之间通过贝叶斯定理相互关联：

$$p(X|Y) = \frac{p(Y|X) \cdot p(X)}{p(Y)}$$

\begin{itemize}
    \item $p(Y|X)$：判别式模型
    \item $p(X)$：生成模型
    \item $p(Y)$：标签先验
    \item $p(X|Y)$：条件生成模型
\end{itemize}

理论上，只要有其中任意两个模型，就能构造出第三个。虽然实践中通常直接训练条件生成模型，但这种关系在扩散模型的 Classifier-Free Guidance 中有重要应用。

\subsection{为什么需要生成模型}

生成模型的核心价值在于处理输出的\textbf{不确定性}和\textbf{多样性}：

\begin{figure}[H]
\centering
\includegraphics[width=0.95\textwidth]{slides-images/slide-014.jpg}
\caption{生成模型的应用场景：语言建模（ChatGPT）、文本生成图像、图像生成视频\protect\footnotemark}
\end{figure}
\footnotetext{Slides 第15页。}

\begin{itemize}
    \item \textbf{语言建模}：输入文本，输出可能有多种合理的回答
    \item \textbf{文本到图像}：一段文字描述可以对应无数张合理的图片
    \item \textbf{图像到视频}：一张静态图片可以衍生出多种可能的未来
\end{itemize}

\begin{knowledgebox}{何时使用生成模型}
只要任务的输出存在\textbf{歧义性}（ambiguity）——即给定输入可能有多种合理的输出——就应该考虑使用生成模型。生成模型建模的是输出的\textbf{分布}，而不是单一的确定性映射。
\end{knowledgebox}

\subsection{生成模型分类体系}

\begin{figure}[H]
\centering
\includegraphics[width=0.95\textwidth]{slides-images/slide-015.jpg}
\caption{生成模型的分类体系：显式密度（可精确/近似计算 $p(X)$） vs 隐式密度（只能采样，不能计算 $p(X)$）\protect\footnotemark}
\end{figure}
\footnotetext{Slides 第16页。}

\begin{importantbox}{生成模型的四大类别}
\begin{enumerate}
    \item \textbf{自回归模型}（Autoregressive）：显式密度，可精确计算 $p(X)$
    \item \textbf{变分自编码器}（VAE）：显式密度，但只能计算 $p(X)$ 的近似下界
    \item \textbf{生成对抗网络}（GAN）：隐式密度，直接采样（单次前向传播）
    \item \textbf{扩散模型}（Diffusion）：隐式密度，迭代采样（需多步推理）
\end{enumerate}
本节课覆盖前三类（自回归、VAE、GAN），下节课讲扩散模型。
\end{importantbox}

\subsection{本章小结}

生成模型学习数据的概率分布 $p(X)$，使得模型能够生成新的数据样本。判别式、生成式、条件生成式三种模型通过贝叶斯定理相互关联。在实践中，条件生成模型最有价值。生成模型的分类体系根据是否能计算密度值以及采样方式分为四大类。

%% ========== Part 3: 最大似然估计 ==========

